% !TEX root = ../double_trouble_main.tex

We first confirm that our method is able to predict well on data simulated from our own model (\cref{eq:proposed_prior}), as a check on our methodology before moving to real data. To that end, we generate 500 data points in each of two populations; we describe the simulation procedure in more detail in~Section S6 of the Supplementary Material \citep{supplement}. So, in each of our 20 folds, the pilot data has size $\vecsamplesize = (25,25)$.

We simulate two different datasets from our model (\cref{eq:proposed_prior}); we choose the hyperparameters to reflect a range of growth rates, sample sizes, and correlations between the two populations in each case. We choose hyperparameters $\allhyper := (\mass, \rate{1},\rate{2}, \corr{1},\corr{2}, \conc{1},\conc{2}) = (100,0.4,0.6,0.5,0.5,1,1)$ for the first dataset and $\allhyper = (1,0.8,0.7,0.5,0.5,1,1)$ for the second dataset. We choose these values for qualitative matches with the observed variant growth in the GnomAD dataset (\cref{fig:pop_gen}) and cancer dataset (\cref{fig:cancer}), respectively. 

The results in \cref{fig:proposed} demonstrate that our method outperforms two potential baselines. Since here all of our data is simulated from our two-population BNP model, we use a single-population BNP model to form our baselines; namely, we compare to the independent and dependent two-population applications of the 3BP approach due to \citet{Masoero2022}. \internal{First, consider the simulation in the upper row of plots.} We see in the upper left plot in \cref{fig:proposed} that our method (solid red line) much more closely matches the behavior of the ground truth (dashed black line) than either the independent (dotted green line) or dependent (dash-dot blue line) applications. As expected, the independent application dramatically overcounts the number of variants in the second follow-up population. Also as expected, the dependent application performs well when the pilot and follow-up proportions are in agreement (at the largest sample size) but misestimates the number of variants elsewhere. By contrast, our method is able to capture the different behaviors of the two populations and account for any shared variants.

 In the lower left corner, the independent application of the 3BP performs the best; its mean line is closest to the ground truth, and its prediction exhibits the smallest variation (shaded range) across folds. Importantly, \internal{we can infer from the lower center plot that} there are extremely few shared variants in \internal{the simulation for the lower row}; see Section S7.3 of the Supplementary Material \citep{supplement} \internal{for raw observed and predicted counts.} That is, \internal{in the simulation for the lower row of plots,} the assumption of the independent application is very close to being correct. So, following Bayesian Occam's Razor \citep[ch.~28 p.~343]{mackaybook}, we expect better performance from the simpler model (the independent application) when its more stringent assumptions are true. That being said, our proposed model still performs well here; its mean line is also very close to the ground truth, by contrast to the dependent application. We conclude that our method is able to handle prediction of the number of new variants under a variety of two-population behaviors --- while the independent application is limited to the special cases where the two populations are known a priori to behave independently.

\begin{figure}[htp]
    \centering
    \includegraphics[width = 0.9\textwidth]{Figs/n3BP_20fold.png}
    \includegraphics[width = 0.9\textwidth]{Figs/n3BP_fast_20fold.png}
    \caption{Predictions when data is sampled from our model (\cref{eq:proposed_prior}). Each row corresponds to a simulated dataset (\cref{sec:simulation}). \emph{Left}: %The dashed black line shows the observed number of variants as a function of the sample number.
    The samples (horizontal axis) appear in the following order: pilot from population 1, pilot from population 2, follow-up from population 1, follow-up from population 2. A vertical dashed gray line separates the two populations in the pilot and follow-up, respectively. Vertical gray shading covers the pilot data. In the follow-up region, we plot mean (lines) and 1 standard deviation intervals (shaded regions) from ground truth (dashed black; also appears for the pilot) and three methods: our proposed method (solid red), the dependent version of the single-population 3BP (dash-dot blue), and the independent version of the 3BP (dotted green). \emph{Center and Right}: For each $\vecoccurrence$ with component values up to 10, we plot the relative residual for predictions from our method (right) and the dependent 3BP (center). \revision{Note: color scales are dataset-specific.} A square is gray when the observed value is strictly less than 2 in at least one fold or $\vecoccurrence = (0,0)$. 
    }
    \label{fig:proposed}
\end{figure}

From the center and right plots in \cref{fig:proposed}, we also see that our method provides better estimates of the numbers of various $\vecoccurrence$-tons. Recall that the independent application predicts zero for anything except private $\vecoccurrence$-tons, and we see that there are in fact non-trivial counts of shared $\vecoccurrence$-tons in the first data set. Given this limitation of the independent application, we focus on plotting results for our method and the dependent application. Darker colors indicate more erroneous predictions, so we conclude in both datasets that our method is providing more accurate predictions than the dependent 3BP application. Note that the observed numbers of $\vecoccurrence$-tons vary dramatically across $\vecoccurrence$ values. So generally we expect more noise in the upper right of our plots of $\vecoccurrence$-tons --- and should therefore take care not to treat every square in these plots equally. We can see the actual observation values for the two datasets in \internal{Fig. S5 (Section~S7.3)} of the Supplementary Material \citep{supplement}; in particular, for the first dataset, the observed number of $\vecoccurrence$-tons varies from an average (across folds) of 1161.2 at $(0,1)$ to 0 at various locations, e.g. (8,3), (10,4), and (7,7) (see the gray boxes in the upper left panel of \internal{Fig. S5 (Section~S7.3)} of the Supplementary Material \citep{supplement}).  We see similar behavior in the real data, as we describe below.

We check that our method is able to handle datasets generated from the (mis-specified) assumptions of the independent and dependent 3BP applications in Sections S7.1 and S7.5 of the Supplementary Material \citep{supplement}; \internal{in particular, in S7.5, we investigate a case where the data is generated from a dependent 3BP where the number of variants grows slowly, and using the independent 3BP for inference results in severe double counting of variants --- but our method continues to work well.}

